\chapter{基本解方法}
\section{广义函数}
\begin{definition}{$\delta$函数}{delta function}
	原点处的 Dirac 分布描述集中在一点的单位总量。若 $a < b$ 且两个端点均非零，对连续函数 $f$ 有
	\begin{equation}
		\label{eq:mmp-dirac-sifting}
		\int_a^bf(x)\delta(x)\d x =
		\begin{cases}
			f(0),&0 \in (a,b)\\
			0,&0 \notin [a,b]
		\end{cases}
	\end{equation}
	因此$f \ast \delta(\xi) = f(\xi)$。
\end{definition}
\begin{theorem}{$\delta$复合函数}{}
	设 $u \in \cont^1(\RR)$ 且在实轴上只有单零点$x_1,x_2,\ldots,x_n$，则
	\begin{align}
		\delta(u(x)) = \sum_{k=1}^n\frac{\delta(x - x_k)}{\abs{u'(x_k)}}.
	\end{align}
\end{theorem}
例如 $a \neq 0$ 时，$\delta(ax) = \delta(x)/\abs a$。

\begin{definition}{线性泛函}{Linear Functional}
	函数空间$\functionspace$到数域$\FF$的映射$\mathcal T$，若$\forall a,b \in \FF,\;u,v \in \functionspace$均有
	\begin{equation}
		\label{eq:mmp-distribution-linearity}
		\mathcal T(au + bv) = a\mathcal Tu + b\mathcal Tv
	\end{equation}
	则 $\mathcal T$ 是 $\functionspace$ 上的线性泛函。全体线性泛函构成代数对偶 $\functionspace^\ast$。给定测试函数空间的拓扑后，连续线性泛函构成连续对偶 $\functionspace'$，下文用 $\inp f\varphi$ 表示泛函 $f$ 在测试函数 $\varphi$ 上的取值。
\end{definition}
\begin{example}{支撑集}{support set}
	记实函数$f$的支撑集
	\begin{align}
		\supp f := \overline{\set x{f(x) \neq 0}}.
	\end{align}
	并记
	\[
		\mathscr L_0(\RR) := \biggset f{\intR\abs{f(x)}\d x < \infty,\;\supp f\,\text{有界}}
	\]
	则$\forall f \in \mathscr L_0(\RR)$，$f$确定了$\cont(\RR)$上一个线性泛函：
	\begin{align}
		\mathcal T_f\bigb{\varphi(x)} \equiv \inp f\varphi := \intR f(x)\varphi(x)\d x \in \RR.
	\end{align}
\end{example}
\begin{theorem}{广义函数的导数}{}
	紧支撑光滑测试函数构成空间
	\begin{align}
		\testspace(\RR) \equiv \cont_{\mathrm c}^\infty(\RR) := \set f{f \in \cont^\infty(\RR),\;\supp f\,\text{有界}}
	\end{align}
	该空间的连续对偶记为 $\testspace'(\RR)$，其元素称为分布（广义函数）。对 $f \in \testspace'(\RR)$，定义其 $n$ 阶导数
	\begin{align}
		\inp{f^{(n)}}\varphi := (-1)^n\inp f{\varphi^{(n)}},\quad\forall\varphi \in \testspace(\RR).
	\end{align}
	特别地， 
	\begin{equation}
		\label{eq:mmp-delta-derivatives}
		\intR\delta^{(n)}(x)\varphi(x)\d x = (-1)^n\varphi^{(n)}(0).
	\end{equation}
	
	泛函意义下，任意一个广义函数都是无穷阶可导的。
\end{theorem}
\begin{theorem}{广义函数的卷积}{}
	设 $f,g \in \testspace'(\RR)$，且至少一个具有紧支撑，则卷积 $f \ast g$ 由下式定义
	\begin{align}
		\inp{f \ast g}\varphi := \Biginp{f(x)}{\biginp{g(y)}{\varphi(x + y)}}.
	\end{align}
	一般地，对常系数微分算子$\DiffOp$
	\begin{equation}
		\label{eq:mmp-convolution-differentiation}
		\DiffOp(f \ast g) = \DiffOp f \ast g = f \ast \DiffOp g.
	\end{equation}
\end{theorem}
\begin{theorem}{广义函数的Fourier变换}{}
	速降函数空间$\Schwartz(\RR) \subset \cont^\infty(\RR)$，%$\forall\varphi\in\Schwartz(\RR)$，有
	\begin{align}
		\Schwartz(\RR) := \set{\varphi \in \cont^\infty(\RR)}{\lim_{\abs x\to\infty}x^m\varphi^{(n)}(x) = 0,\quad m,n = 0,1,\ldots}.
	\end{align}
	比如$\varphi(x) = p(x)\e^{-ax^2}$，$p(x)$是多项式，$a > 0$，则$\varphi(x) \in \Schwartz(\RR)$。
	
	速降测试函数的连续对偶 $\Schwartz'(\RR)$ 称为温和分布空间。给定 $f \in \Schwartz'(\RR)$，其 Fourier 变换及逆变换定义为
	\begin{align}
		\begin{aligned}
			\inp{\FT[f]}\varphi &= \inp f{\FT[\varphi]}\\
			\inp{\FT\inv[f]}\varphi &= \inp f{\FT\inv[\varphi]}
		\end{aligned},\quad\forall\varphi \in \Schwartz.
	\end{align}
	Fourier变换作用转移到测试函数上，它们保持着经典意义下的基本性质。
\end{theorem}
\begin{theorem}{广义函数序列的收敛性}{}
	设广义函数列$\{f_n\}$及广义函数$f$，若对测试函数空间的$\varphi$都有
	\begin{equation}
		\label{eq:mmp-distributional-convergence}
		\lim_{n\to\infty}\inp{f_n}\varphi = \inp f\varphi.
	\end{equation}
	则称 $f_n$ 在分布意义下弱收敛到 $f$，记作 $f_n\rightharpoonup f$。
\end{theorem}
\begin{theorem}{}{}
	若 $f_n,f \in \Schwartz'(\RR)$，则求导保持这种收敛性
	\begin{equation}
		\label{eq:mmp-weak-derivative-continuity}
		f_n\rightharpoonup f \implies f_n'\rightharpoonup f'.
	\end{equation}
\end{theorem}
\begin{theorem}{高维广义函数}{}
	多重指标 $\alpha = (\alpha_1,\ldots,\alpha_n)$ 的分量为非负整数，记 $\p^\alpha = \p_{x_1}^{\alpha_1}\cdots\p_{x_n}^{\alpha_n}$。与一维类似，函数空间
	\begin{subequations}
		\label{eq:mmp-multidimensional-test-spaces}
		\begin{gather}
			\testspace(\RR^n) \equiv \cont_{\mathrm c}^\infty(\RR^n) := \set f{f \in \cont^\infty(\RR^n),\;\supp f\,\text{有界}},\\
			\Schwartz(\RR^n) := \Bigset{\varphi \in \cont^\infty(\RR^n)}{\lim_{\norm X\to\infty}\norm X^m\p^\alpha\varphi(X) = 0,\quad m = 0,1,\ldots,\;\forall\alpha}.
		\end{gather}
	\end{subequations}
	局部可积函数 $f$ 在测试函数空间上给出连续线性泛函
	\begin{equation}
		\label{eq:mmp-regular-multidimensional-distribution}
		\inp f\varphi = \int_{\RR^n}f(X)\varphi(X)\d X,\quad\forall\varphi \in \testspace(\RR^n).
	\end{equation}
	
	取 $\alpha = (k_1,\ldots,k_n)$，$k = k_1 + \cdots + k_n$。广义函数 $f \in \testspace'(\RR^n)$ 的偏导数定义为
	\begin{equation}
		\label{eq:mmp-distribution-multiindex-derivative}
		\inp{\p^\alpha f}{\varphi} := (-1)^k\inp{f}{\p^\alpha\varphi},\quad\varphi \in \testspace(\RR^n).
	\end{equation}
	特别地，$f = \delta$时
	\begin{equation}
		\label{eq:mmp-delta-multiindex}
		\inp{\p^\alpha\delta}{\varphi} = (-1)^k(\p^\alpha\varphi)(\bm 0).
	\end{equation}
\end{theorem}
\begin{example}{$\delta$函数与Laplace算子}{}
	\begin{equation}
		\delta(x_1,\ldots,x_n) = \begin{cases}
			\frac1{2\pi}\laplacian\ln r,&n = 2\\
			-\frac1{(n - 2)S_{n}}\laplacian\frac1{r^{n-2}},&n \geqslant 3
		\end{cases}
	\end{equation}
	其中$r = \norm{\bm x}$，$S_n$为$n$维单位球的表面积。特别地，$S_3 = 4\pi$。
\end{example}
若至少一个分布具有紧支撑，高维卷积同样由配对定义
\begin{equation}
	\label{eq:mmp-multidimensional-convolution}
	\inp{f \ast g}\varphi := \Biginp{f(X)}{\biginp{g(Y)}{\varphi(X + Y)}}.
\end{equation}

\section{\texorpdfstring{$\PDEop u = 0$}{Pu=0}型方程的基本解} % {$\$}
讨论用基本解方法求解方程
\[
	\PDEop u(M) = f(M),\quad M \in \RR^n,
\]
其中，$\PDEop$是常系数线性偏微分算子。

在测试函数上按分布意义满足$\PDEop u = f$的解称为\textbf{广义解}。分布可以逐阶求导。若$u$由足够光滑的函数给出且$f$连续，则此等式也逐点成立，$u$是经典解。
\begin{definition}{$\PDEop u = 0$型方程的基本解}{Basic Solution of Pu=0}
	方程
	\begin{align}
		\PDEop U(M) = \delta(M)
	\end{align}
	的解$U(M)$称作$\PDEop u = 0$型方程的基本解。
\end{definition}
对一般的函数$f$，其对应的解是一般的源所产生的物理场。故基本解也叫点源函数。

设$U$是基本解，源$f$为光滑紧支撑函数。若$U$局部可积，可写
\begin{align}
	u(M) := U \ast f(M) = \int U(M - N)f(N)\d N.
\end{align}
由积分叠加原理
\[
	\PDEop u(M) = \int\PDEop U(M - N)f(N)\d N = \int\delta(M - N)f(N)\d N = f(M).
\]
即$u$满足方程$\PDEop u = f$。更一般地，只要分布卷积存在，就有$\PDEop(U \ast f) = (\PDEop U) \ast f = f$；例如$f$具有紧支撑时成立。再加任意齐次解仍满足同一非齐次方程。
\begin{example}{}{}
	求下面方程的基本解：
	\[
		y' + ay = f(x),\quad a > 0
	\]
	
	基本解$U$满足
	\[
		U' + aU = \delta(x),
	\]
	取在$x < 0$处为零的因果基本解，由
	\begin{equation}
		\label{eq:mmp-first-order-causal-kernel}
		\dd x(U\e^{ax}) = \delta(x)\e^{ax} = \delta(x) \implies U(x) = \Heaviside(x)\e^{-ax}.
	\end{equation}
	因此，只要下式积分收敛，便得到原方程的一个特解
	\begin{equation}
		\label{eq:mmp-causal-first-order-solution}
		y(x) = U \ast f(x) = \intR\Heaviside(\xi)\e^{-a\xi}f(x - \xi)\d\xi = \int_{-\infty}^xf(\xi)\e^{-a(x-\xi)}\d\xi.
	\end{equation}
\end{example}
\begin{example}{3维Helmholtz方程}{}
	求3维Helmholtz方程的基本解：
	\[
		\laplacian u + cu = 0
	\]
	
	由于只有一个点源，方程具有对称性，可设方程有球对称基本解$U(r)$，
	当$r > 0$时，$\delta(r) = 0$，
	\[
		\frac1{r^2}\dd r\Bigp{r^2\dv Ur} + cU = \frac1r\Bigb{\dd[2]r(rU) + crU} = 0,
	\]
	
	当$c = 0$时，化为Laplace方程
	\[
		(rU)'' = 0 \implies U(r) = \frac Ar + \cancel{B}.
	\]
	记$B_r$为半径为$r$的球内部分，则
	\[
		\int_{B_r}\laplacian U\d V = \int_{B_r}\delta(x,y,z)\d V = 1.
	\]
	由Green公式
	\[
		\int_{B_r}\laplacian U\d V = \oint_{\p B_r}\pv Un\d S = -\frac A{r^2} \cdot 4\pi r^2 = 1 \implies U(r) = -\frac1{4\pi r}.
	\]
	
	当$c < 0$时，记$k := \sqrt{-c}$
	\[
		(rU)'' - k^2rU = 0 \implies U = A\frac{\e^{-kr}}r + B\frac{\e^{kr}}r.
	\]
	由积分条件
	\[
		\int_{B_r}\laplacian U - k^2U\d V = \oint_{\p B_r}\pv Un\d S - k^2\int_{B_r}U\d V = 1,
	\]
	且
	\begin{align*}
		&\quad\oint_{\p B_r}\Bigp{\pp\rho\frac{\e^{k\rho}}\rho}_rr^2\sin\theta\d\theta\nd\phi - k^2\int_{B_r}\frac{\e^{k\rho}}\rho\rho^2\sin\theta\d\rho\nd\theta\nd\phi\\
		&= 4\pi(kr - 1)\e^{kr} - 4\pi\bigb{(k\rho - 1)\e^{k\rho}}_0^r = -4\pi.
	\end{align*}
	故$-4\pi(A + B) = 1$，即
	\begin{equation}
		\label{eq:mmp-modified-helmholtz-fundamental}
		U(r) = -\frac{\e^{-kr}}{4\pi r},\; -\frac{\e^{kr}}{4\pi r}.
	\end{equation}
	
	当$c > 0$时，记$k := \sqrt c$
	\[
		U = A\frac{\cos kr}r + B\frac{\sin kr}r.
	\]
	$\sin kr/r$在原点光滑，是齐次解。点源归一化给出$-4\pi A = 1$，而$B$可任取。例如取$B = 0$得到
	\begin{equation}
		\label{eq:mmp-real-helmholtz-fundamental}
		U(r) = -\frac{\cos kr}{4\pi r}.
	\end{equation}
\end{example}
\section{Poisson方程Green函数法}
首先引入Green公式。
\begin{theorem}{Green公式}{Green Formula}
	设非空有界开集$\Omega \subset \RR^n(n \geqslant 2)$满足边界$\p\Omega$光滑，则
	
	设$u,v$在$\bar\Omega$的一个开邻域上二阶连续可微，有
	\begin{gather}
		\label{Green's 1st formula}
		\int_\Omega v\laplacian u\d V = \oint_{\p\Omega}v\pv un\d S - \int_\Omega\grad v \cdot \grad u\d V;\\
		\label{Green's 2nd formula}
		\int_\Omega(u\laplacian v - v\laplacian u)\d V = \oint_{\p\Omega}\biggp{u\pv vn - v\pv un}\d S.
	\end{gather}
	%其中$\p u/\p n=\grad u\cdot\uvec{n}$。
\end{theorem}
\begin{proof}
	由散度定理
	\begin{align*}
		\int_\Omega\div\bm A\d V = \oint_{\p\Omega}\bm A \cdot \uvec{n}\d S,
	\end{align*}
	取$\bm A = v\grad u$，
	可得式\eqref{Green's 1st formula}：
	\[
		\oint_{\p\Omega}v\pv un\d S \equiv \oint_{\p\Omega}v\grad u \cdot \uvec{n}\d S = \int_\Omega\div(v\grad u)\d V = \int_\Omega(\grad v \cdot \grad u + v\laplacian u)\d V.
	\]
	将$u,v$互换位置并与原式相减，即得式\eqref{Green's 2nd formula}。
\end{proof}
\paragraph{Poisson方程第I边值问题}
\begin{align}
	\begin{cases}
		\laplacian u = -f(M),&M \in V \subseteq \RR^3\\
		\eval u_{\p V} = \varphi(M),
	\end{cases}
\end{align}
在介电常数为$\varepsilon$的均匀介质中，电势满足$\laplacian u = -\rho/\varepsilon$。这里的源项对应电荷体密度$\rho = \varepsilon f$，边界电势为$\varphi$。

由叠加原理，$u = v + w$，$v,w$分别满足
\begin{align*}
	\begin{cases}
		\laplacian v = -f(M),\\
		\eval v_{\p V} = 0,
	\end{cases}\quad
	\begin{cases}
		\laplacian w = 0,\\
		\eval w_{\p V} = \varphi(M),
	\end{cases}
\end{align*}
其中$v$是边界接地时体内电荷产生的电势，$w$是体内无电荷时由给定边界电势产生的解。
\begin{definition}{Poisson第I边值问题的Green函数}{}
	定解问题
	\begin{align}
		\begin{cases}
			\laplacian G(M;N) = -\delta(M - N),\\
			\eval G_{\p V} = 0.
		\end{cases}
	\end{align}
	的解$G(M;N)$称为Poisson第I边值问题的Green函数。
\end{definition}
在介电常数为$\varepsilon$的均匀介质中，若边界接地，置于$N$点的点电荷$q$在$M$点产生的电势为$qG(M;N)/\varepsilon$。
\begin{align}
	u(M) = \int_Vf(N)G(M;N)\d N - \oint_{\p V}\varphi(N)\pv Gn\d S.
\end{align}
\begin{proof}
	\begin{align*}
		u(M) &= \int_Vu(N)\delta(M - N)\d N = -\int_Vu(N)\,\laplacian G(M;N)\d N\\
		&= -\int_V u\,\laplacian G\d N + \int_VG\,\laplacian u\d N - \int_VG\,\underline{\laplacian u}\d N\tag{Green II}\\
		&= \oint_{\p V}\cancel{G}\pv un - \underline{u}\pv Gn\d S + \int_VG f\d N\\
		&= -\oint_{\p V}\varphi\pv Gn\d S + \int_VGf\d N.
	\end{align*}
\end{proof}
Fourier方法是求Green函数的基本方法，但对于一些特殊的区域，可以采用一些特殊方法，如镜像法。
\begin{method}{镜像法}{}
	分解$G = U + g$，$U$满足$\laplacian U = -\delta(M - N)$，可取
	\begin{equation}
		\label{eq:mmp-free-laplace-kernels}
		U = \begin{cases}
			-\frac1{2\pi}\ln r,&n = 2\\[1ex]
			\frac1{4\pi r},&n = 3
		\end{cases}
	\end{equation}
	%其中$r$是$M,N$间距离。
	而$g$满足 
	\[
		\begin{cases}
			\laplacian g = 0,\\
			\eval g_{\p V} = -U,
		\end{cases}
	\]
	是$N$点的电荷在边界上诱导的电荷所产生的电势。
	
	区域外的点源在$V$内产生的电势满足Laplace方程。若区域\emph{外}的虚设电荷产生的电势满足$g$的边界条件，由解的唯一性，它在$V$内就等于$g$。这种方法叫镜像法。% 它的关键困难在于如何在区域外合适地虚设电荷，对应某些特殊的区域如半空间、球域等等，可以用较直观的方法找到。
\end{method}
\begin{example}{镜像法应用·半空间}{}
	求上半空间Poisson方程第I边值问题的Green函数。
	
	$V$内$N(\xi,\eta,\zeta)$点的正电荷$\varepsilon$在空间$(x,y,z)$产生的电势为
	\[
		U_0 = \frac1{4\pi r} = \frac1{4\pi\sqrt{(x - \xi)^2 + (y - \eta)^2 + (z - \zeta)^2}},
	\]
	可虚设电荷$-\varepsilon$于$N$关于$z = 0$平面对称的点$M_1(\xi,\eta,-\zeta)$，产生的电势为
	\[
		U_1 = \frac{-1}{4\pi\sqrt{(x - \xi)^2 + (y - \eta)^2 + (z + \zeta)^2}},
	\]
	在边界上$\eval{U_0}_{z=0} = -\eval{U_1}_{z=0}$
    \begin{equation}
        \label{eq:mmp-half-space-dirichlet-green}
        \begin{aligned}
            G = \frac1{4\pi}\bigg[&\frac1{\sqrt{(x - \xi)^2 + (y - \eta)^2 + (z - \zeta)^2}} -\\
            &\qquad\frac1{\sqrt{(x - \xi)^2 + (y - \eta)^2 + (z + \zeta)^2}}\bigg],
        \end{aligned}
    \end{equation}
	边界方向导数
	\[
		\biggeval{\pv Gn}_{\zeta=0} = \biggeval{-\pv G\zeta}_{\zeta=0} = \frac{-z}{2\pi[(x - \xi)^2 + (y - \eta)^2 + z^2]^{3/2}}
	\]
	故对于上半空间Dirichlet问题
	\[
		\begin{cases}
		\laplacian u = 0,&z > 0\\
		\eval u_{z=0} = \varphi(x,y)
	\end{cases}
	\]
	解的Poisson公式
	\begin{equation}
		\label{eq:mmp-half-space-poisson-integral}
		u(x,y,z) = \frac z{2\pi}\int_{\RR^2}\frac{\varphi(\xi,\eta)\d\xi\nd\eta}{[(x - \xi)^2 + (y - \eta)^2 + z^2]^{3/2}}
	\end{equation}
	\tcblower
	二维情况
	\begin{subequations}
		\label{eq:mmp-half-plane-dirichlet-green}
		\begin{align}
			G &= \frac1{4\pi}\ln\frac{(x - \xi)^2 + (y + \eta)^2}{(x - \xi)^2 + (y - \eta)^2},\\
			u(x,y) &= \frac y\pi\intR\frac{\varphi(\xi)\d\xi}{(x - \xi)^2 + y^2}
		\end{align}
	\end{subequations}
\end{example}

\begin{example}{镜像法应用·球域}{}
	设球半径为$R$，场点$M$和源点$N$到球心的距离分别为$r$和$\rho$，两条径向矢量的夹角为$\psi$。球面的Poisson核等于Dirichlet Green函数对源点的外法向导数取负号：
	\begin{gather}
		\begin{cases}
			\laplacian_N G(M;N) = -\delta(M - N),&0 \leqslant \rho < R\\
			\eval G_{\rho=R} = 0
		\end{cases}\notag\\
		\implies -\biggeval{\pv Gn}_{\rho=R} = \frac{R^2 - r^2}{4\pi R(R^2 + r^2 - 2Rr\cos\psi)^{3/2}}.\label{eq:mmp-ball-poisson-kernel}
	\end{gather}
	故球内Dirichlet问题，解的Poisson形式
	\begin{gather}
		\begin{cases}
			\laplacian u = 0,&0 \leqslant r < R\\
			\eval u_{r=R} = \varphi(\theta,\phi)
		\end{cases}\notag\\
		\implies u(r,\theta,\phi) = \frac1{4\pi R}\oint_{S_R}\frac{(R^2 - r^2)\varphi(\theta_0,\phi_0)\d S_0}{(R^2 + r^2 - 2Rr\cos\psi)^{3/2}}.\label{eq:mmp-ball-poisson-integral}
	\end{gather}
\end{example}

\subparagraph*{Fourier法}
Fourier法是求Green函数的基本方法，主要思想是按照特征函数作广义Fourier展开，包括分离变量与积分变换。
\begin{example}{Fourier法应用}{}	
	求矩形区域$\Omega = [0,L] \times [0,M]$上第I类边值Poisson方程的Green函数
	\begin{align*}
		\begin{cases}
			\laplacian G(x,y;\xi,\eta) = -\delta(x - \xi,y - \eta),\quad x,\xi \in [0,L];\;y,\eta \in [0,M]\\
			G(0,y) = G(L,y) = G(x,0) = G(x,M) = 0,
		\end{cases}
	\end{align*}
	有特征值 
	\[
		G(x,y) = \sum_{m,n}C_{mn}\sin\frac{m\pi x}L\sin\frac{n\pi y}M.
	\]
	系数 
	{\footnotesize
	\begin{align*}
		\Bigb{\Bigp{\frac{m\pi}L}^2 + \Bigp{\frac{n\pi}M}^2}C_{mn} &= \frac4{LM}\int_0^M\int_0^L\delta(x - \xi,y - \eta)\sin\frac{m\pi x}L\sin\frac{n\pi y}M\d x\nd y\\
		&= \frac4{LM}\sin\frac{m\pi\xi}L\sin\frac{n\pi\eta}M.
	\end{align*}}
	后略。
\end{example}
\begin{example}{Fourier法应用·有限高}{}
	\begin{equation*}
		\begin{cases}
			\laplacian u = 0,\quad x > 0,y \in [0,a]\\
			u(x,0) = \varphi(x),\enspace u(x,a) = \psi(x),\\
			u(0,y) = 0,
		\end{cases}
	\end{equation*}
	取光滑且紧支撑于$(0,\infty)$的边界数据，并取$x \to \infty$时衰减的解。有
	\begin{align*}
		u(x,y) &= -\biggb{\intRplus\varphi(\xi)\biggeval{\pv Gn}_{\eta=0}\d\xi + \intRplus\psi(\xi)\biggeval{\pv Gn}_{\eta=a}\d\xi}\\
		&= \intRplus\varphi(\xi)\biggeval{\pv G\eta}_{\eta=0}\d\xi - \intRplus\psi(\xi)\biggeval{\pv G\eta}_{\eta=a}\d\xi
	\end{align*}
	Green函数在$\eta$和$y$方向都采用正弦基，对源坐标$\eta$求导后得
	\begin{align*}
		\pv G\eta &= \sum_{n=1}^\infty4n\intRplus\frac{\sin\omega\xi\sin\omega x}{n^2\pi^2 + a^2\omega^2}\d\omega \cdot \cos\frac{n\pi\eta}a\sin\frac{n\pi y}a;\\
		u(x,y) &= \sum_{n=1}^\infty4n\intRplus\intRplus\frac{\sin\omega\xi\sin\omega x}{n^2\pi^2 + a^2\omega^2}[\varphi(\xi) - (-1)^n\psi(\xi)]\d\omega\nd\xi \cdot \sin\frac{n\pi y}a.
	\end{align*}
	最后一式用于$0 < y < a$，两端的边界值由内部极限取得。
\end{example}
\paragraph{$^\ast$Poisson方程第II, III边值问题}
\begin{align}
	\begin{cases}
		\laplacian u = -f(M),&M \in V \subseteq \RR^n\\
		\Bigp{\alpha u + \beta\pv un}_{\p V} = \varphi(M),&\alpha,\beta \neq 0.
	\end{cases}
\end{align}
设$V$为连通有界区域，取常系数Robin条件，且齐次问题只有零解（例如$\alpha/\beta > 0$）。相应的Green函数满足
\begin{equation}
	\label{eq:mmp-robin-green-problem}
	\begin{cases}
		\laplacian G = -\delta(M - N),&M,N \in V \subseteq \RR^n\\
		\Bigp{\alpha G + \beta\pv Gn}_{\p V} = 0,&\alpha,\beta \neq 0.
	\end{cases}
\end{equation}
原问题的解
\begin{align}
	u(M) &= \int_Vf(N)G(M;N)\d N + \frac1\beta\oint_{\p V}\varphi(N)G(M;N)\d S\\
	&= \int_Vf(N)G(M;N)\d N - \frac1\alpha\oint_{\p V}\varphi(N)\pv Gn\d S
\end{align}
$\beta = 0$即第I边值问题。$\alpha = 0$时，记给定的外法向导数为$g = \varphi/\beta$。对第II边值问题，常数是齐次解。若同时要求$\laplacian G = -\delta$与齐次Neumann边界，对全域积分便与散度定理矛盾。因此可选用以下任一修正
\begin{equation}
	\label{eq:mmp-normalized-neumann-kernels}
	\begin{cases}
		\laplacian G = -\delta(M - N) + \frac1v,\\
		\biggeval{\pv Gn}_{\p V} = 0,
	\end{cases}\enspace\text{或}\quad
	\begin{cases}
		\laplacian G = -\delta(M - N),\\
		\biggeval{\pv Gn}_{\p V} = -\frac1s,
	\end{cases}
\end{equation}
其中$v,s$为$V$的体积和表面积。有解
\begin{align}
	u(M) = \int_Vf(N)G(M;N)\d N + \oint_{\p V}g(N)G(M;N)\d S + \const.
\end{align}
相容性条件
\begin{align}
	\int_Vf(M)\d M + \oint_{\p V}g(M)\d S = 0.
\end{align}
\section{初值问题的基本解方法}
本节主要用基本解方法来求解发展方程，如$u_t = \PDEop u$型方程
\begin{align}
	\begin{cases}
		u_t = \PDEop u + f(M,t),&M \in \RR^n,t > 0\\
		u(M,0) = \varphi(M).
	\end{cases}
\end{align}
这里所涉及的$\PDEop$是关于空间变量$M$的常系数线性偏微分算子。

\begin{definition}{$u_t = \PDEop u$型方程初值问题的基本解}{Basic Solution of ut=Pu}
	基本解$U(M,t)$满足
	\begin{align}
		\begin{cases}
			U_t = \PDEop U,&M \in \RR^n,t > 0\\
			U(M,0) = \delta(M).
		\end{cases}
	\end{align}
\end{definition}
有
\begin{align}
	u(M,t) = U(M,t) \ast \varphi(M) + \int_0^tU(M,t - \tau) \ast f(M,\tau)\d\tau.
\end{align}
证明过程略。

以及$u_{tt} = \PDEop u$型方程
\begin{align}
	\begin{cases}
		u_{tt} = \PDEop u + f(M,t),&M \in \RR^n,t > 0\\
		u(M,0) = \varphi(M),\enspace u_t(M,0) = \psi(M).
	\end{cases}
\end{align}
\begin{definition}{$u_{tt} = \PDEop u$型方程初值问题的基本解}{Basic Solution of utt=Pu}
	基本解$U(M,t)$满足
	\begin{align}
		\begin{cases}
			U_{tt} = \PDEop U,&M \in \RR^n,t > 0\\
			U(M,0) = 0,\enspace U_t(M,0) = \delta(M).
		\end{cases}
	\end{align}
\end{definition}
有
\begin{align}
	\begin{aligned}
		u(M,t) =&\;U \ast \psi + \pp t(U \ast \varphi) + \int_0^tU(M,t - \tau) \ast f(M,\tau)\d\tau.
	\end{aligned}
\end{align}
\paragraph{*混合问题的Green函数}对于发展方程的混合问题通常用分离变量法或积分变换法求解，当然也可以用Green函数（点源函数）法。
\begin{example}{一维波动方程的混合问题}{}
	一维波动方程的混合问题
	\begin{align*}
		\begin{cases}
			u_{tt} = a^2u_{xx} + f(x,t),\quad x \in (0,L),\;t > 0\\
			u(0,t) = u(L,t) = 0\\
			u(x,0) = \varphi(x),\quad u_t(x,0) = \psi(x)
		\end{cases}
	\end{align*}
	Green函数$G(x,t;\xi)$满足 
	\begin{align*}
		\begin{cases}
			G_{tt} = a^2G_{xx},\quad x,\xi \in (0,L),\;t > 0\\
			G(0,t) = G(L,t) = 0\\
			G(x,0) = 0,\quad G_t(x,0) = \delta(x - \xi)
		\end{cases}
	\end{align*}
	利用Fourier方法可得
	\begin{equation}
		\label{eq:mmp-mixed-string-retarded-green}
		G(x,t;\xi) = \sum_{n=1}^\infty\frac2{n\pi a}\sin\frac{n\pi\xi}L\sin\frac{n\pi x}L\sin\frac{n\pi at}L,
	\end{equation}
	进而解
	\begin{align}
		\begin{aligned}
			u(x,t) = {}&\int_0^L\psi(\xi)G\d\xi + \pp t\int_0^L\varphi(\xi)G\d\xi +\\
			&\int_0^t\int_0^L f(\xi,\tau)G(x,t - \tau;\xi)\d\xi\nd\tau
		\end{aligned}
	\end{align}
\end{example}

% \input{Append.tex}
